% ==============================================================================
\chapter{Conclusion \& Future Perspectives}
\label{ch:conclusion}
% ==============================================================================

\section{Summary of Achievements}
This thesis has detailed the systematic development of a real-time, client-side interactive Digital Twin platform utilizing Isogeometric Analysis (IGA) and projection-based Reduced Order Modeling (ROM) techniques. The project successfully bridged advanced computational continuum mechanics with modern web software engineering.

The primary milestones completed during the development include:
\begin{enumerate}
    \item \textbf{Theoretical Foundation}: Established the structural baseline kernels to handle smooth NURBS evaluations and verified baseline projection errors via Proper Orthogonal Decomposition.
    \item \textbf{Full Order Model Construction}: Built a native 2D tensor-product physical solver environment capable of generating accurate displacement behaviors under open uniform knot distributions.
    \item \textbf{Nonlinear Kinematics Integration}: Extended the element integration loops to compute the Green-Lagrange strain tensor, introducing an incremental residual-driven Newton-Raphson loop to trace geometrically nonlinear deformations.
    \item \textbf{Hyper-Reduction Optimization}: Successfully deployed the Energy-Conserving Sampling and Weighting (ECSW) hyper-reduction technique, breaking the classical lifting bottleneck and preserving topological column alignment across shape updates.
    \item \textbf{Digital Twin Deployment}: Built an asynchronous, client-side interactive application layer utilizing JSON-serialized ROM payloads that achieves stable convergence in under 16ms, ensuring a smooth 60FPS user interaction budget.
\end{enumerate}

\section{Future Perspectives and Research Paths}
While the platform successfully meets its immediate pedagogical and performance objectives, several optimization vectors are identified for subsequent research:

\subsection{Expansion to 3D Continuum Formulations}
The existing physics core operates on 2D plane stress and plane strain mechanics. Scaling the tensor-product loop to process trivariate NURBS solids would enable the simulation of complex 3D volumes (e.g., thick turbine components or composite structures), significantly broadening the platform's industrial applicability.

\subsection{Multi-Patch Topology Assembly}
The current geometric mapping restricts the structural mesh domain to a single, continuous bivariate patch topology. Integrating multi-patch blending mechanics via G-continuity or penalty constraints (such as Nitsche's method or Mortar element formulations) would allow the web platform to represent intricate, non-isomorphic multipatch CAD definitions exactly.

\subsection{Adaptive Online Hyper-Reduction Sampling}
Currently, the ECSW integration weights are computed entirely during the offline training pass. If the user shifts the parameter sliders far beyond the sampled training bounds, approximation accuracy decreases. Integrating on-the-fly basis enrichment techniques (such as localized kriging interpolation or automated online snapshot sampling) would allow the twin to retain structural accuracy dynamically.


